\documentclass[12pt]{article}
\def\activityid{F09-X-v3}\usepackage[letterpaper,margin=.65in,top=.60in,bottom=.65in]{geometry}
\usepackage[T1]{fontenc}
\usepackage[default]{sourcesanspro}
\usepackage{microtype,tikz,array,tabularx,fancyhdr,amsmath,amssymb}
\usepackage[hidelinks]{hyperref}
\usetikzlibrary{calc,arrows.meta}
\setlength{\parindent}{0pt}\setlength{\parskip}{7pt}
\setlength{\headheight}{0pt}\setlength{\footskip}{26pt}\setlength{\emergencystretch}{2em}
\pagestyle{fancy}\fancyhf{}\renewcommand{\headrulewidth}{0pt}\renewcommand{\footrulewidth}{.4pt}
\fancyfoot[L]{\scriptsize Bellingham Math Circle / Week 9 / \activityid}
\fancyfoot[R]{\scriptsize \thepage}
\definecolor{ink}{gray}{.12}\definecolor{soft}{gray}{.95}\color{ink}
\newcommand{\PageTitle}[2]{{\fontsize{10}{12}\selectfont\bfseries WEEK 9\quad BOUNCE LAB\hfill #1\par}\vspace{3pt}{\fontsize{26}{29}\selectfont\bfseries #2\par}\vspace{2pt}}
\newcommand{\Subhead}[1]{\par\vspace{3pt}{\large\bfseries #1}\par}
\newcommand{\RuleBox}[1]{\par\begingroup\setlength{\fboxsep}{8pt}\colorbox{soft}{\parbox{\dimexpr\linewidth-16pt\relax}{#1}}\endgroup\par}
\newcommand{\WriteBox}[2]{\par\textbf{#1}\par\vspace{2pt}\begin{tikzpicture}\draw[black!40,line width=.5pt] (0,0) rectangle (\dimexpr\linewidth-.5pt\relax,#2);\end{tikzpicture}\par}
\newcommand{\RookBoard}[3]{\begin{tikzpicture}[x=#3,y=#3]\draw[step=1,black!45] (0,0) grid (#1,#2);\draw[thick] (0,0) rectangle (#1,#2);\node[font=\Large] at ({#1-.5},.5){$\star$};\draw[-{Latex},thick] (0,-.35)--(#1,-.35);\draw[-{Latex},thick] ({#1+.35},#2)--({#1+.35},0);\end{tikzpicture}}
\newcommand{\Table}[3]{\begin{tikzpicture}[x=#3,y=#3]\draw[step=1,black!25] (0,0) grid (#1,#2);\draw[thick] (0,0) rectangle (#1,#2);\fill (0,0)circle(3pt);\draw[-{Latex},very thick] (0,0)--(.7,.7);\node[below] at ({#1/2},-.1){width #1};\node[rotate=90,above] at (-.2,{#2/2}){height #2};\end{tikzpicture}}
\newcommand{\GraphNodes}{\coordinate(A)at(0,0);\coordinate(B)at(3,0);\coordinate(C)at(1.5,2.6);\coordinate(D)at(1.5,.95);}
\newcommand{\Kfour}[1]{\begin{tikzpicture}[scale=#1]\GraphNodes\draw[very thick](A)--(B)--(C)--(A)--(D)--(B) (D)--(C);\foreach \n in {A,B,C,D}{\node[circle,draw,fill=white,minimum size=22pt,font=\bfseries]at(\n){\n};}\end{tikzpicture}}
\newcommand{\House}[1]{\begin{tikzpicture}[scale=#1]\coordinate(A)at(0,0);\coordinate(B)at(3,0);\coordinate(C)at(3,2);\coordinate(D)at(0,2);\coordinate(E)at(1.5,3.3);\draw[very thick](A)--(B)--(C)--(D)--(A) (D)--(E)--(C);\foreach \n in {A,B,C,D,E}{\node[circle,draw,fill=white,minimum size=22pt,font=\bfseries]at(\n){\n};}\end{tikzpicture}}




\newcommand{\StudentPage}[1]{{\fontsize{10}{12}\selectfont\bfseries Week 9 / Billiards lab / #1\par}\vspace{10pt}}
\newcommand{\Problem}[2]{\par\noindent\textbf{Problem #1:} #2\par}
\newcommand{\Space}[1]{\par\begin{tikzpicture}\draw[black!35] (0,0) rectangle (\dimexpr\linewidth-.5pt\relax,#1);\end{tikzpicture}\par}
\newcommand{\Piles}[1]{\begin{center}\begin{tikzpicture}[x=1in,y=1in]\foreach \x in {0,...,#1}{\draw[black!50,rounded corners] ({\x*2.05},0) rectangle ({\x*2.05+1.8},1.3);}\end{tikzpicture}\end{center}}

\newcommand{\Rooms}[5]{\begin{tikzpicture}[x=#5,y=#5]\draw[step=1,black!20](0,0)grid(#3,#4);\foreach\x in{0,#1,...,#3}{\draw[thick](\x,0)--(\x,#4);}\foreach\y in{0,#2,...,#4}{\draw[thick](0,\y)--(#3,\y);}\fill(0,0)circle(3pt);\end{tikzpicture}}
\newcommand{\PlainGrid}[3]{\begin{tikzpicture}[x=#3,y=#3]\draw[step=1,black!30](0,0)grid(#1,#2);\fill(0,0)circle(3pt);\end{tikzpicture}}

\begin{document}

\StudentPage{Extra / Grades 6--7}
\Problem{1}{Use a unit square, unfolded into the square grids below. Each little square is one whole reflected table. Start at the black dot. Slope means vertical distance divided by horizontal distance. Draw slopes $1/2$, $2/3$, and $3/2$, one on each grid. Stop at the first grid vertex after the starting dot. Mark each wall crossing before it.}
\begin{center}\PlainGrid{4}{3}{.45in}\hspace{.6in}\PlainGrid{4}{3}{.45in}\end{center}
\begin{center}\PlainGrid{3}{4}{.45in}\end{center}
\Problem{2}{Record each first vertex as (horizontal distance, vertical distance). Fold the whole-room counts back to predict the corner. Count the wall crossings before it.}
\begin{center}\renewcommand{\arraystretch}{1.9}\begin{tabular}{c|p{1.35in}|p{1.4in}|p{1.3in}}Slope&First vertex&Corner&Bounces\\\hline$1/2$&&&\\$2/3$&&&\\$3/2$&&&\end{tabular}\end{center}
\newpage
\StudentPage{Extra / Grades 6--7}
\Problem{3}{On the left grid draw the line of slope $2/4$. On the right draw the line of slope $3/4$. Stop each at its first grid vertex. Record the vertex, corner, and bounces; compare $2/4$ with your $1/2$ path.}
\begin{center}\PlainGrid{4}{4}{.52in}\hspace{.7in}\PlainGrid{4}{4}{.52in}\end{center}
\Space{.7in}
\Problem{4}{Use positive fractions $p/q$ with no common factor between $p$ and $q$. Find a slope whose path ends top-right after four bounces. Find one that ends bottom-right after three bounces. Record each slope and its first unfolded vertex, then check the wall counts.}
\Space{1.0in}
\Problem{5}{Now use slope $\sqrt2$. If its unfolded path from (0,0) passed through a positive integer point (u,v), its slope would be $v/u$. Record what this would say about $\sqrt2$. You may use the fact that $\sqrt2$ is not a ratio of integers.}
\Space{.8in}
\newpage
\StudentPage{Extra / Grades 6--7}
\Problem{6}{Start at the midpoint of the left wall of this unit square, heading up-right with slope 1. Trace four bounces. Continue for four more. Mark a point inside the square that your path never reaches.}
\begin{center}\begin{tikzpicture}[x=3in,y=3in]\draw[thick](0,0)rectangle(1,1);\draw[dashed,black!25](0,.5)--(1,.5) (.5,0)--(.5,1);\fill(0,.5)circle(3pt);\draw[-{Latex},thick](0,.5)--(.18,.68);\end{tikzpicture}\end{center}
\Problem{7}{Compare your midpoint path with the slope-$\sqrt2$ path from (0,0). For each, record whether it can reach a corner. For the midpoint path, draw a small region that the path never enters.}
\Space{1in}
\Problem{8}{A known theorem says the slope-$\sqrt2$ path from (0,0) comes arbitrarily close to every point in the square. This is called dense. Circle which evidence below establishes density: ``I did not hit a corner'' / ``the stated theorem.'' Use your midpoint path to give an example showing why the other claim is insufficient.}
\Space{.6in}

\end{document}
